\documentclass[11pt]{article}

\usepackage[margin=1in]{geometry}
\usepackage{amsmath,amssymb,amsthm}
\usepackage{enumitem}

\newtheorem{theorem}{Theorem}
\newtheorem{corollary}{Corollary}
\newtheorem{remark}{Remark}

\newcommand{\Sym}{\operatorname{Sym}}
\newcommand{\Stab}{\operatorname{Stab}}
\newcommand{\Aut}{\operatorname{Aut}}
\newcommand{\Cay}{\operatorname{Cay}}
\newcommand{\Cent}{\operatorname{C}}

\title{A Stronger Rigidity Formulation for the Zaks Cycle}
\author{}
\date{}

\begin{document}
\maketitle

\section*{Main point}

The current theorem about the Zaks Hamiltonian cycle can be reframed as a sharp
case of a more general rigidity principle for cyclic orderings of $S_n$.

The general principle does not depend on the pancake graph. It only uses the
left regular action of $S_n$ and the elementary fact that the centralizer of an
$n$-cycle in $S_n$ is the cyclic group generated by it.

\section*{General cyclic-order rigidity theorem}

Let $C$ be a cyclic ordering of the elements of $S_n$. Equivalently, $C$ is a
Hamiltonian cycle on the complete graph with vertex set $S_n$; no Cayley graph
structure is needed here.

For $g \in S_n$, write $L_g$ for left multiplication:
\[
  L_g(\sigma)=g\sigma.
\]

\begin{theorem}[Cyclic-order rigidity for $S_n$]
Let $n\geq 3$. Let $C$ be a cyclic ordering of the elements of $S_n$, and let
$c \in S_n$ be an $n$-cycle. Assume that $C$ is invariant under the left
translation $L_c$:
\[
  L_c(C)=C.
\]
Then
\[
  \Stab_{L(S_n)}(C)
  \subseteq
  \{\,L_g : g c g^{-1}\in\{c,c^{-1}\}\,\}.
\]
In particular,
\[
  |\Stab_{L(S_n)}(C)| \leq 2n.
\]
\end{theorem}

\begin{proof}
Since $C$ is a cyclic ordering, its abstract automorphism group is dihedral. The
translation $L_c$ acts on $C$ as a rotation of order $n$. If $L_g$ stabilizes
$C$, then inside the automorphism group of the cycle it conjugates this rotation
either to itself or to its inverse:
\[
  L_g L_c L_g^{-1}=L_c
  \quad\text{or}\quad
  L_g L_c L_g^{-1}=L_c^{-1}.
\]
The left regular action is faithful, hence
\[
  g c g^{-1}=c
  \quad\text{or}\quad
  g c g^{-1}=c^{-1}.
\]
It remains to count the possible $g$'s. The centralizer of an $n$-cycle in
$S_n$ is exactly
\[
  \Cent_{S_n}(c)=\langle c\rangle.
\]
Thus there are $n$ elements with $g c g^{-1}=c$. The elements with
$g c g^{-1}=c^{-1}$ form one coset of this centralizer, hence there are another
$n$ of them. Therefore the stabilizer is contained in a set of size $2n$.
\end{proof}

The containing group is the dihedral group of the $n$-gon whose cyclic order is
defined by $c$. If
\[
  c=(a_1\,a_2\,\cdots\,a_n),
\]
then it is the group of symmetries of the polygon
\[
  a_1,a_2,\ldots,a_n.
\]

\section*{Sharpness: the Zaks cycle attains the bound}

Now return to the pancake graph
\[
  P_n=\Cay(S_n,\{r_2,\ldots,r_n\}),
\]
with generators acting by right multiplication. Let $Z_n$ be the Zaks
Hamiltonian cycle, and set
\[
  \rho=r_{n-1}r_n.
\]

The Zaks recursion gives
\[
  L_\rho(Z_n)=Z_n.
\]
Thus the general theorem applies with $c=\rho$.

Moreover, the palindromic property of the Zaks word gives the reflection
\[
  L_{r_n}(Z_n)=Z_n.
\]
Since
\[
  r_n \rho r_n=\rho^{-1},
\]
we get the full dihedral subgroup
\[
  \langle L_\rho,L_{r_n}\rangle \cong D_n
\]
inside the left-regular stabilizer.

The general rigidity theorem gives the reverse inclusion, because no
left-regular stabilizer of a cyclic ordering invariant under $L_\rho$ can have
more than $2n$ elements. Therefore
\[
  \Stab_{L(S_n)}(Z_n)
  =
  \langle L_\rho,L_{r_n}\rangle
  \cong D_n.
\]

\begin{corollary}[Sharpness]
The bound
\[
  |\Stab_{L(S_n)}(C)|\leq 2n
\]
is sharp. It is attained by the Zaks Hamiltonian cycle in the pancake graph.
\end{corollary}

\section*{Minimal Coxeter interpretation}

This is the Coxeter/Weyl dictionary, in only the amount needed for the paper.

The group $S_n$ is the Weyl group of type $A_{n-1}$. Its simple Coxeter
generators are the adjacent transpositions
\[
  s_i=(i\ i+1).
\]
The pancake generator $r_k$ is not a simple reflection. Instead,
\[
  r_k
  =
  \text{the longest element of the parabolic subgroup }
  \langle s_1,\ldots,s_{k-1}\rangle
  \cong S_k.
\]
In particular,
\[
  r_n=w_0,
\]
the longest element of the Weyl group $S_n$.

Also,
\[
  \rho=r_{n-1}r_n
\]
is an $n$-cycle, hence a Coxeter element of type $A_{n-1}$. Thus the stabilizer
found above is not an arbitrary dihedral group:
\[
  \langle \rho,r_n\rangle
  =
  \langle c,w_0\rangle,
\]
the natural dihedral group generated by a Coxeter element and the longest
element.

\begin{remark}
For readers not interested in Coxeter theory, the whole argument can be read
purely in terms of $S_n$: $n$-cycles give the natural rotations, and the
centralizer of an $n$-cycle is the cyclic group it generates. The Coxeter
language explains why the Zaks choices $\rho=r_{n-1}r_n$ and $r_n$ are
geometrically distinguished.
\end{remark}

\section*{Effect on the paper}

This reframing changes the structure of the proof.

\begin{enumerate}[label=\arabic*.]
  \item Prove the general cyclic-order rigidity theorem for $S_n$.
  \item Prove that the Zaks recursion gives $L_\rho(Z_n)=Z_n$.
  \item Prove that the palindrome gives $L_{r_n}(Z_n)=Z_n$.
  \item Conclude exactness from the general rigidity theorem.
  \item Use Deng--Zhang to pass from the left-regular stabilizer to the full
  graph automorphism stabilizer for $n\geq 5$.
\end{enumerate}

So the Zaks theorem becomes: a general upper bound for cyclic orderings of
$S_n$, plus a natural pancake-graph construction that attains the bound.

\end{document}

